La structure de donnée permettant d'obtenir une complexité en $O(m + n\log n)$ est celle de la file e priorité minimale. On peut l'implémenter par un tas binaire minimal, ce qui permet d'obtenir une complexité de $O(\log n)$ pour les opérations d'insertion et de suppression du minimum. Cependant, il existe une structure de données plus efficace appelée le \href{https://www.google.com/url?sa=t&source=web&rct=j&opi=89978449&url=https://lirias.kuleuven.be/retrieve/353036&ved=2ahUKEwjO2rrPwpKTAxW7dqQEHWRbFtkQFnoECDcQAQ&usg=AOvVaw3Ylzd0cMfgLtrKYd2Y9Ppi}{tas de Fibonacci}, qui permet d'obtenir une complexité de $O(1)$ pour l'insertion et la diminution de clé, et de $O(\log n)$ pour la suppression du minimum. En utilisant l'une de ces implémentations, on peut donc obtenir une complexité totale de $O(m + n\log n)$ pour l'algorithme de Dijkstra, ou bien pour $A^*$, car la seule différence entre les deux algorithmes est l'utilisation d'une heuristique pour guider la recherche, ce qui a lieu dans l'ajout d'un élément dans la file de priorité, et qui ne change pas la complexité de l'algorithme dans le cas général.


